\chapter{Firmware Developments}
\label{ch:dev}

The changes described here were made in the SR5E1 CAN library and, where
the SDK had no hook, in weakly-defined Motor Control functions. Each
subsection states the defect, the design constraint, and the resulting
code. Listings are shortened; complete functions are in
Appendix~\ref{ch:app}.

\section{Mutually exclusive \texttt{EXECUTE\_COMMAND}}
\label{sec:exec}

\subsection{The defect}
\texttt{EXECUTE\_COMMAND} (identifier \texttt{0x33}) carries a bit-mask of
machine-level actions. The original decoder was a sequence of independent
\texttt{if} blocks: one for start, one for stop, one for reset, and so on.
Two consequences followed.

\begin{itemize}[leftmargin=1.6em]
  \item Because start was coded as
        \texttt{if (RxArr[0] \& EXE\_START\_MOTOR)} and stop as
        \texttt{else if (!(RxArr[0] \& EXE\_START\_MOTOR))}, \emph{every}
        frame produced either a start or a stop. There was no way to send
        a reset, a ping or a fault acknowledgement without also commanding
        the machine to stop. The \texttt{else if} was not a second flag; it
        was the complement of the first.
  \item Reset lived in its own \texttt{if} after the start/stop path, so a
        single CAN frame could emit two FCP frames---a stop and a reset---
        back to back. On a bus already running at the diagnostic period of
        \SI{100}{\milli\second}, the extra burst was visible as a glitch in
        machine state and as a pair of ACKs that a careless consumer could
        mis-attribute.
\end{itemize}

\subsection{The repair}
The decoder was rewritten around an explicit \texttt{validCommand} flag
and a chain of \texttt{if / else if} that commits to at most one Motor
Control command per CAN frame. The start bit still means start; its
absence no longer means stop. Stop, reset, fault acknowledgement, encoder
alignment and ping are each a distinct flag, and only the first one that
matches is translated. If no recognised flag is set,
\texttt{validCommand} stays false and nothing is fed to the FCP layer.

\begin{lstlisting}[style=cstyle,caption={Decoded EXECUTE\_COMMAND path (simplified).},label=lst:exec]
case EXECUTE_COMMAND:
    bool validCommand = true;
    msg[0] = MC_PROTOCOL_CODE_EXECUTE_CMD;
    msg[1] = 0x01;

    if (RxArr[0] & EXE_START_MOTOR) {
        msg[2] = MC_PROTOCOL_CMD_START_MOTOR;
    } else if (RxArr[0] & EXE_STOP_MOTOR) {
        msg[2] = MC_PROTOCOL_CMD_STOP_MOTOR;
    } else if (RxArr[0] & EXE_RESET) {
        msg[2] = MC_PROTOCOL_CMD_RESET;
    } else if (RxArr[0] & EXE_FAULT_ACK) {
        msg[2] = MC_PROTOCOL_CMD_FAULT_ACK;
    } else if (RxArr[0] & EXE_ENCODER_ALIGN) {
        msg[2] = MC_PROTOCOL_CMD_ENCODER_ALIGN;
    } else {
        validCommand = false;
    }

    if (validCommand) {
        msg[3] = SelfCRCCalc(msg);
        for (count = 0; count < (msg[1] + 3); count++) {
            (void) CFCP_RX_IRQ_Handler(&pFDCAN, msg[count]);
        }
    }
    break;
\end{lstlisting}

The listing is the version that was kept. An intermediate attempt, visible
on the project slides, still treated the absence of the start bit as stop
and then appended a second FCP feed for reset. That version removed the
double start/stop but not the double transmission; the \texttt{validCommand}
form removes both.

\subsection{Flags consumed by the decoder}
\Cref{tab:exec-flags} records the bit-mask as it appears in the DBC and in
the firmware. Bits not listed are reserved and must be zero; the decoder
ignores them so that a future GUI revision can claim them without breaking
old boards.

\begin{table}[ht]
\centering
\caption{EXECUTE\_COMMAND bit-mask.}
\label{tab:exec-flags}
\begin{tabular}{ccl}
\toprule
Byte & Bit & Flag \\
\midrule
0 & 1 & \texttt{EXE\_START\_MOTOR} \\
0 & 2 & \texttt{STOP\_RAMP} \\
0 & 3 & \texttt{RESET} \\
0 & 5 & \texttt{EXE\_PING} \\
0 & 6 & \texttt{FAULT\_ACK} \\
0 & 7 & \texttt{ENCODER\_ALIGN} \\
1 & 0 & \texttt{EXE\_IQDREF\_CLEAR} \\
1 & 1 & \texttt{GET\_BOARD\_INFO} \\
\bottomrule
\end{tabular}
\end{table}

\section{Independent $I_d$ and $I_q$ references}
\label{sec:id-iq}

\subsection{Why the pair is not enough}
\texttt{SET\_CURRENT\_ID\_IQ} (identifier \texttt{0x3A}) already writes both
axes in one frame. Characterisation, however, almost always holds one axis
and moves the other: a flux-map identification holds $I_q = 0$ and sweeps
$I_d$; a torque-linearity test holds $I_d = 0$ (or a weakening value) and
sweeps $I_q$. Doing that with a paired command requires the bench to
remember the axis it does not want to change and to rewrite it on every
frame. Independent commands remove that state from the bench.

\subsection{Protocol codes}
Two Motor Control codes that the original library did not use were claimed:

\begin{lstlisting}[style=cstyle,caption={New protocol codes for independent current references.}]
#define MC_PROTOCOL_CODE_SET_CURRENT_ID_REF  0x0F
#define MC_PROTOCOL_CODE_SET_CURRENT_IQ_REF  0x11
\end{lstlisting}

On the CAN side the corresponding identifiers, after the offset rule and
after collision avoidance with the ramp window, are \texttt{0x3B}
(\texttt{SET\_CURRENT\_ID}) and \texttt{0x3C} (\texttt{SET\_CURRENT\_IQ}).
Each carries a 16-bit signed little-endian set-point; the FCP length byte
is therefore \texttt{0x02}.

\begin{lstlisting}[style=cstyle,caption={DYNO2DW cases for independent $I_d$ / $I_q$.},label=lst:id-iq-dyno]
case SET_CURRENT_ID:
    msg[0] = MC_PROTOCOL_CODE_SET_CURRENT_ID_REF;
    msg[1] = 0x02;
    msg[2] = RxArr[0];
    msg[3] = RxArr[1];
    msg[4] = SelfCRCCalc(msg);
    break;

case SET_CURRENT_IQ:
    msg[0] = MC_PROTOCOL_CODE_SET_CURRENT_IQ_REF;
    msg[1] = 0x02;
    msg[2] = RxArr[0];
    msg[3] = RxArr[1];
    msg[4] = SelfCRCCalc(msg);
    break;
\end{lstlisting}

\subsection{UI application}
On the SDK side the codes are applied with the existing setters, after a
16-bit signed reconstruction that matches the Cortex little-endian load of
an \texttt{int16\_t}:

\begin{lstlisting}[style=cstyle,caption={UI dispatch for independent current references.}]
case MC_PROTOCOL_CODE_SET_CURRENT_ID_REF: {
    int16_t hIdRef = (int16_t)buffer[0]
                   + ((int16_t)buffer[1] << 8);
    UI_SetIdRef(&pHandle->_Super, hIdRef);
    bNoError = true;
} break;

case MC_PROTOCOL_CODE_SET_CURRENT_IQ_REF: {
    int16_t hIqRef = (int16_t)buffer[0]
                   + ((int16_t)buffer[1] << 8);
    UI_SetIqRef(&pHandle->_Super, hIqRef);
    bNoError = true;
} break;
\end{lstlisting}

No range check is performed at this layer. The Motor Control current
regulators already saturate against the configured maximum; duplicating
that saturation on the CAN boundary would have created two sources of
truth. A NACK is returned only if the drive handle is invalid.

\section{Torque ramp ($I_q$)}
\label{sec:iq-ramp}

The torque ramp was already present in the SDK as
\texttt{MC\_PROTOCOL\_CODE\_SET\_TORQUE\_RAMP} /
\texttt{UI\_ExecTorqueRamp}. What the original library lacked was the
DYNO2DW case that forwards a CAN frame onto that code. The payload is
six bytes: a 32-bit signed final $I_q$ and a 16-bit unsigned duration in
milliseconds.

\begin{lstlisting}[style=cstyle,caption={Torque-ramp translation and SDK dispatch.}]
case SET_IQ_RAMP:
    msg[0] = MC_PROTOCOL_CODE_SET_TORQUE_RAMP;
    msg[1] = 0x06;
    /* RxArr[0..3] = final Iq, RxArr[4..5] = duration */
    ...
    break;

case MC_PROTOCOL_CODE_SET_TORQUE_RAMP: {
    int32_t  torque   =  (int32_t)buffer[0]
                      + ((int32_t)buffer[1] <<  8)
                      + ((int32_t)buffer[2] << 16)
                      + ((int32_t)buffer[3] << 24);
    uint16_t duration =  (uint16_t)buffer[4]
                      + ((uint16_t)buffer[5] << 8);
    bNoError = UI_ExecTorqueRamp(&pHandle->_Super, torque, duration);
} break;
\end{lstlisting}

A duration of zero means ``step'': the reference jumps in the next control
period. Non-zero durations are converted, inside the speed/torque
controller, into a number of ticks at \texttt{STCFrequencyHz} and into an
increment per tick, exactly as the flux ramp does in the next section.

\section{Flux ramp ($I_d$)}
\label{sec:id-ramp}

\subsection{A missing primitive}
Unlike the torque ramp, a flux ramp did not exist. The SDK could set a
constant $I_d$ and it could ramp $I_q$; it could not ramp $I_d$ over a
specified time. That gap is visible in any flux-weakening characterisation
and in any test that must enter or leave $I_d < 0$ without a step in
magnetising current.

The implementation therefore adds a protocol code,
\texttt{MC\_PROTOCOL\_CODE\_SET\_FLUX\_RAMP}, a DYNO2DW case
\texttt{SET\_ID\_RAMP} of length 6, and three new (weak) functions that
mirror the torque-ramp call chain:

\begin{center}
\texttt{UI\_ExecFluxRamp} $\rightarrow$
\texttt{MCI\_ExecFluxRamp} $\rightarrow$
\texttt{STC\_ExecCurrentRamps}
\end{center}

\subsection{UI and MCI layers}
The UI function is a thin selected-drive indirection. The MCI function
records the request on the motor-control interface handle and marks it as
not yet executed, so that the regular control-period task, rather than the
CAN ISR, performs the numerical work.

\begin{lstlisting}[style=cstyle,caption={UI and MCI entry points of the flux ramp.},label=lst:flux-ui]
__weak bool UI_ExecFluxRamp(UI_Handle_t *pHandle,
                            qd_t Vqd, uint16_t hDurationms)
{
    MCI_Handle_t *pMCI = pHandle->pMCI[pHandle->bSelectedDrive];
    MCI_ExecFluxRamp(pMCI, Vqd, hDurationms);
    return true;
}

__weak void MCI_ExecFluxRamp(MCI_Handle_t *pHandle,
                             qd_t Iqdref, uint16_t hDurationms)
{
    pHandle->lastCommand   = MCI_EXECFLUXTRAMP;
    pHandle->hFinalFlux    = Iqdref.d;
    pHandle->hDurationms   = hDurationms;
    pHandle->CommandState  = MCI_COMMAND_NOT_ALREADY_EXECUTED;
}
\end{lstlisting}

\subsection{STC ramp law}
The speed/torque controller stores flux in Q16 (the reference is
multiplied by $2^{16}$). Given a current reference $I_d(0)$, a final value
$I_d^\star$ and a duration $T$ in milliseconds, the number of remaining
steps and the per-step increment are

\begin{align}
N &= \left\lfloor T \cdot f_{\mathrm{STC}} / 1000 \right\rfloor + 1, \\
\Delta &= \frac{(I_d^\star - I_d(0)) \cdot 2^{16}}{N}.
\end{align}

If $T = 0$, the increment is zero and the Q16 reference is loaded
immediately. The controller is switched to torque mode
(\texttt{STC\_TORQUE\_MODE}) so that the outer speed loop cannot fight the
current ramps. The same function updates the $I_q$ ramp when called from
the torque path; the two axes share the duration-to-ticks conversion but
keep separate remaining-step counters and separate increments
(\texttt{RampRemainingStepId}, \texttt{IncDecAmountId}).

\begin{lstlisting}[style=cstyle,caption={STC flux-ramp update (excerpt).},label=lst:stc]
__weak bool STC_ExecCurrentRamps(SpeednTorqCtrl_Handle_t *pHandle,
                                 int16_t hFinalFlux,
                                 uint16_t hDurationms)
{
    int16_t hCurrentReferenceId = STC_GetFluxRef(pHandle);
    STC_SetControlMode(pHandle, STC_TORQUE_MODE);

    if (hDurationms == 0u) {
        pHandle->FluxRef            = (int32_t)hFinalFlux * 65536;
        pHandle->RampRemainingStepId = 0u;
        pHandle->IncDecAmountId      = 0;
    } else {
        pHandle->TargetFinalId = hFinalFlux;
        uint32_t wAux = (uint32_t)hDurationms
                      * (uint32_t)pHandle->STCFrequencyHz;
        wAux /= 1000u;
        pHandle->RampRemainingStepId = wAux + 1u;
        int32_t wAux1 = ((int32_t)hFinalFlux
                       - (int32_t)hCurrentReferenceId) * 65536;
        wAux1 /= (int32_t)pHandle->RampRemainingStepId;
        pHandle->IncDecAmountId = wAux1;
    }
    return true;
}
\end{lstlisting}

The DYNO2DW case that feeds this chain preserves the $I_q$ component of
the current \texttt{qd\_t} pair---it reads \texttt{MCI\_GetIqdref},
overwrites only \texttt{.d}, and then calls the ramp---so that commanding
an $I_d$ ramp cannot accidentally zero $I_q$.

\section{Packed PI register access}
\label{sec:pi-pack}

\subsection{The representation}
The MCSDK does not store a floating-point $K_p$. It stores an integer
numerator and a power-of-two divisor~\cite{astromhagglund2006,stmcsdk}:

\begin{equation}
K_p^{\mathrm{eff}} = \frac{K_p}{2^{n}}, \qquad
K_i^{\mathrm{eff}} = \frac{K_i}{2^{m}}.
\end{equation}

\texttt{PID\_SetKP} writes $K_p$; \texttt{PID\_SetKPDivisorPOW2} writes
$n$. A CAN consumer that can set only one of the two cannot reproduce a
tuned loop, and a consumer that sets them in two messages can be
preempted by a control period that sees a mixed pair.

\subsection{The packing}
A single 32-bit \texttt{wValue} therefore carries both halves:

\begin{center}
\begin{tabular}{cc}
\toprule
Bits 31--16 & Bits 15--0 \\
\midrule
divisor $n$ (unsigned) & gain $K_p$ (unsigned, used as signed 16-bit) \\
\texttt{(wValue >> 16) \& 0xFFFF} & \texttt{wValue \& 0xFFFF} \\
\texttt{PID\_SetKPDivisorPOW2()} & \texttt{PID\_SetKP()} \\
\bottomrule
\end{tabular}
\end{center}

The same layout is used for $K_i$ and for the three loops (speed, flux,
torque). Reading is the inverse operation, so a \texttt{GET\_REG} of
\texttt{MC\_PROTOCOL\_REG\_SPEED\_KP} round-trips.

\begin{lstlisting}[style=cstyle,caption={Packing and unpacking of SPEED\_KP.},label=lst:pi]
/* UI_SetReg */
case MC_PROTOCOL_REG_SPEED_KP: {
    uint16_t kp_gain     = (uint16_t)(wValue & 0xFFFF);
    uint16_t kp_div_pow2 = (uint16_t)(((uint32_t)wValue >> 16) & 0xFFFF);
    PID_SetKP(pMCT->pPIDSpeed, (int16_t)kp_gain);
    PID_SetKPDivisorPOW2(pMCT->pPIDSpeed, kp_div_pow2);
} break;

/* UI_GetReg */
case MC_PROTOCOL_REG_SPEED_KP: {
    bRetVal = ((int32_t)(PID_GetKPDivisor(pMCT->pPIDSpeed) & 0xFFFFu) << 16)
            |  ((uint16_t)(PID_GetKP(pMCT->pPIDSpeed) & 0xFFFFu));
} break;
\end{lstlisting}

The DBC exposes the two halves as two signals on the same message
(\texttt{REGISTER\_VALUE}, \texttt{REGISTER\_VALUE\_DENOM}), so the GUI of
\cref{ch:gui} can show ``Num / Den POW2'' without the operator doing the
shifts. The firmware never sees those signal names; it only sees
\texttt{wValue}.

\section{Separating ACK/NACK from \texttt{GET\_REG}}
\label{sec:ack}

\subsection{The collision}
Every well-formed command produces a response. A write produces a one-byte
acknowledgement, \texttt{0xF0} for success and \texttt{0xFF} for failure.
A read produces a four-byte register value. In the original library both
responses were emitted under the same CAN identifier. A consumer that had
just sent \texttt{GET\_REG} could not know whether the next frame was the
value it asked for or an ACK from a previous write that had been delayed
by bus arbitration. A consumer that had just sent \texttt{SET\_REG} could
mistake a leftover \texttt{REG\_VALUE} for an ACK.

\subsection{Dispatch on frame size}
The FCP TX handler already knew the size of the frame it was serialising.
That size is a reliable discriminant: acknowledgements are not four bytes,
and register values are. The handler therefore sets the identifier to
\texttt{0x10} by default and overrides it to \texttt{0x11} when
\texttt{TxFrame.Size == 4}.

\begin{lstlisting}[style=cstyle,caption={Identifier selection in CFCP\_TX\_IRQ\_Handler.},label=lst:ack]
__weak void CFCP_TX_IRQ_Handler(CFCP_Handle_t *pHandle)
{
    if (FCP_TRANSFER_IDLE != pBaseHandle->TxFrameState) {
        TxHeader.Identifier = CAN_TX_MSG_ID_START;   /* 0x10 */
        if (pBaseHandle->TxFrame.Size == 4) {
            TxHeader.Identifier = 17U;               /* 0x11 */
        }
        /* build tx_data[], submit to FDCAN */
    }
}
\end{lstlisting}

On the DBC side the same split is expressed as two messages,
\texttt{ACK\_NACK} and \texttt{REG\_VALUE}. A further DBC defect---treating
\texttt{0xF0} and \texttt{0xFF} themselves as identifiers---is repaired in
\cref{ch:dbc}; it is mentioned here because it was a symptom of the same
confusion between ``the code that means ACK'' and ``the identifier that
carries it''.

\section{Summary of firmware changes}
\Cref{tab:dev-sum} collects the identifiers, codes and payload sizes that
the rest of the thesis treats as the contract.

\begin{table}[ht]
\centering
\caption{Firmware contract after the developments of this chapter.}
\label{tab:dev-sum}
\begin{tabular}{llll}
\toprule
CAN ID & Direction & FCP code / register & Payload \\
\midrule
\texttt{0x10} & MCU $\rightarrow$ dyno & ACK/NACK & 1 byte (\texttt{0xF0}/\texttt{0xFF}) \\
\texttt{0x11} & MCU $\rightarrow$ dyno & \texttt{REG\_VALUE} & 4 bytes \\
\texttt{0x31} & dyno $\rightarrow$ MCU & \texttt{SET\_REG} & id + packed value \\
\texttt{0x32} & dyno $\rightarrow$ MCU & \texttt{GET\_REG} & register id \\
\texttt{0x33} & dyno $\rightarrow$ MCU & \texttt{EXECUTE\_CMD} & bit-mask \\
\texttt{0x37} & dyno $\rightarrow$ MCU & speed ramp & 6 bytes \\
\texttt{0x3B} & dyno $\rightarrow$ MCU & $I_d$ ref & 2 bytes \\
\texttt{0x3C} & dyno $\rightarrow$ MCU & $I_q$ ref & 2 bytes \\
\texttt{0x3E} & dyno $\rightarrow$ MCU & $I_d$ ramp & 6 bytes \\
\texttt{0x3F} & dyno $\rightarrow$ MCU & $I_q$ ramp & 6 bytes \\
\texttt{0x40}+ & MCU $\rightarrow$ dyno & diagnostics & 4 bytes, period \SI{100}{\milli\second} \\
\bottomrule
\end{tabular}
\end{table}
